\section{Đánh giá và kết luận}
\label{sec:danh-gia-ket-luan}

Trong các chương trước, đồ án đã tìm hiểu về lý thuyết nhóm, các cách ép mô hình học máy tuân theo tính invariance và equivariance, cũng như các chứng minh về độ phức tạp mẫu và kết quả chạy thực nghiệm. Việc thêm tính đối xứng vào mạng nơ-ron đem lại nhiều lợi ích. Tuy nhiên, để đánh giá khách quan, ta cần xem xét kỹ hơn: Liệu có phải lúc nào cũng cần dùng geometric machine learning? Phương pháp này có giới hạn gì, và sắp tới lĩnh vực này sẽ phát triển theo hướng nào? Phần cuối của đồ án sẽ phân tích các vấn đề này.

\subsection{Thảo luận: Có nên luôn dùng mô hình equivariant?}
\label{subsec:thao-luan}

Thấy được thành công của NequIP \cite{batzner2022nequip} khi giảm được rất nhiều dữ liệu cần thiết để dự đoán năng lượng phân tử, hay các chứng minh toán học giúp giảm độ phức tạp mẫu \cite{tahmasebi2023exact}, ta dễ nghĩ rằng cứ dùng đối xứng là mô hình sẽ tốt. Tuy nhiên, theo đánh giá từ tutorial NeurIPS 2025 \cite{jegelka2025tutorial}, thực tế không đơn giản như vậy.

\subsubsection{Xét về quy mô dữ liệu và scaling laws}

Lý do phổ biến nhất để dùng inductive bias là giúp tiết kiệm dữ liệu. Điều này rất quan trọng trong hóa học lượng tử hay thiết kế thuốc, vì tính toán mô phỏng tốn rất nhiều thời gian và tiền bạc. Khi có ít dữ liệu, việc đưa thẳng quy luật không gian vào mô hình giúp nó hội tụ nhanh hơn là để nó tự học từ đầu.

Tuy nhiên, với các foundation models hiện nay được train trên hàng tỷ dữ liệu, nhiều người cho rằng chỉ cần ném dữ liệu vào mạng transformer thì nó sẽ tự học được các tính đối xứng ngầm mà không cần thiết kế kiến trúc phức tạp từ trước. 

Dù vậy, các quan sát thực nghiệm trong \cite{jegelka2025tutorial} cho thấy: ngay cả khi có nhiều dữ liệu và máy mạnh, nếu xét cùng một mức chi phí huấn luyện, mô hình equivariant vẫn có sai số thấp hơn mô hình thường. Nghĩa là việc dùng tính đối xứng giúp cải thiện trực tiếp định luật tỷ lệ, làm quá trình học hiệu quả hơn ở mọi quy mô.

\begin{figure}[H]
    \centering
    \begin{tikzpicture}
        \begin{axis}[
            width=10cm, height=6.2cm,
            xlabel={Chi phí tính toán huấn luyện},
            ylabel={Loss tập kiểm tra},
            xmin=1, xmax=10, ymin=0, ymax=1.1,
            xtick=\empty, ytick=\empty,
            legend pos=north east,
            legend cell align=left,
            axis lines=left,
            domain=1:10, samples=80,
        ]
            \addplot[thick, red!70!black] {0.9/(x^0.45)};
            \addlegendentry{Mô hình thông thường}
            \addplot[thick, blue!60!black] {0.7/(x^0.55)};
            \addlegendentry{Mô hình equivariant}
        \end{axis}
    \end{tikzpicture}
    \caption[Tác động của tính đối xứng lên scaling laws]{Mô hình equivariant thường có hiệu suất tốt hơn mô hình thông thường ở cùng một mức chi phí huấn luyện. (Vẽ lại dựa trên nhận xét từ \cite{jegelka2025tutorial})}
    \label{fig:scaling-law-equivariant}
\end{figure}

\subsubsection{Loại đối xứng trong dữ liệu}

Một vấn đề khác là liệu tính đối xứng có thực sự đúng với mọi dữ liệu hay không. Thực tế có hai trường hợp:

\begin{itemize}
    \item \textit{Đối xứng chính xác:} Đây là các quy luật vật lý bắt buộc. Ví dụ, năng lượng phân tử không đổi khi ta xoay nó đi. Lúc này, dùng invariance là hoàn toàn chính xác.
    \item \textit{Đối xứng xấp xỉ:} Nhiều dữ liệu chỉ có đối xứng tương đối. Ví dụ trong xử lý ảnh, con mèo ở góc nào thì vẫn là con mèo, nhưng vị trí của nó có thể mang thông tin về bối cảnh ảnh. Nếu ép mô hình phải invariance cứng nhắc với phép tịnh tiến, ta có thể vô tình làm mất thông tin vị trí hữu ích.
\end{itemize}

Ngoài ra, theo \cite{jegelka2025tutorial}, một số tập dữ liệu phân tử (benchmark) nổi tiếng đã bị tác giả âm thầm canonicalization sẵn trước khi công bố. 

\begin{warningbox}{Lỗi khi dùng benchmark đã căn lề sẵn}
Nếu dữ liệu phân tử đã bị căn lề, các mẫu không còn xoay ngẫu nhiên nữa. Khi đó, dùng mô hình equivariant (vốn cố tình bỏ qua thông tin góc xoay) lại làm mô hình học kém hơn mạng MLP bình thường, vì ta vô tình vứt bỏ đi thông tin góc xoay đã được cố định sẵn. Do đó, cần kiểm tra kỹ tập dữ liệu trước khi chọn kiến trúc.
\end{warningbox}

\subsubsection{Tính đa năng so với chuyên môn hóa}

Cuối cùng là việc chọn giữa làm một bài toán cụ thể hay làm một mô hình đa năng. Một mô hình equivariant thiết kế cho nhóm xoay $SO(3)$ dự đoán phân tử rất tốt, nhưng không dùng được cho bài toán khác không có không gian 3D. Ngược lại, các foundation models dù học chậm ban đầu nhưng lại xử lý được nhiều loại đối xứng khác nhau tùy theo ngữ cảnh dữ liệu. Do đó, việc chọn geometric machine learning phụ thuộc vào mục đích làm công cụ chuyên biệt hay làm mô hình nền tảng.

\subsection{Tổng hợp ưu điểm và hạn chế}
\label{subsec:uu-nhuoc-diem}

Từ các phần trên, đồ án đúc kết lại các ưu nhược điểm của geometric machine learning:

\textbf{Về ưu điểm:} Lợi ích lớn nhất là học nhanh và cần ít dữ liệu hơn (nhờ inductive bias). Thay vì cần hàng ngàn mẫu ảnh để mô hình biết ảnh xoay ngược vẫn là ảnh đ", ta cài đặt sẵn quy luật đó vào mạng. Việc này cũng giúp mô hình dự đoán đúng trên các dữ liệu bị biến đổi góc độ chưa từng gặp lúc train. 

Bên cạnh đó, nhờ dùng chung trọng số, số tham số cần học ít đi nên giảm được hiện tượng overfitting. Cuối cùng, như đã chỉ ra ở Chương 3, dùng các kỹ thuật như group averaging không làm mất khả năng học được mọi hàm số của mạng gốc.

\textbf{Về hạn chế:} Khó khăn nhất là việc lập trình và thiết kế đòi hỏi kiến thức toán nặng về group theory và representation theory. Hơn nữa, việc tính toán để duy trì tính đối xứng cho các nhóm lớn (như nhóm hoán vị) làm tốn rất nhiều thời gian chạy máy. 

Thêm vào đó, mô hình equivariant rất khó áp dụng cho các bài toán sinh dữ liệu, vì nhiệm vụ này yêu cầu phá vỡ tính đối xứng để tạo ra cấu trúc mới từ đầu vào đối xứng. Cuối cùng, nếu đoán sai nhóm đối xứng hoặc dùng dữ liệu đã bị căn lề sẵn như đã nói ở trên, hiệu suất mô hình sẽ giảm mạnh.

\begin{table}[H]
    \centering
    \small
    \renewcommand{\arraystretch}{1.3}
    \begin{tabular}{p{0.46\textwidth}p{0.46\textwidth}}
        \toprule
        \textbf{Ưu điểm} & \textbf{Hạn chế} \\
        \midrule
        Tiết kiệm dữ liệu và học nhanh hơn vì không phải học lại các phép biến đổi không gian. & Lập trình khó, cần nền tảng toán học sâu để thiết kế và sửa lỗi. \\
        Dự đoán ổn định ngay cả khi dữ liệu đầu vào bị biến đổi. & Chi phí tính toán cao khi áp dụng group averaging trên các nhóm lớn. \\
        Giảm lượng tham số nhờ dùng chung trọng số. & Khó làm các bài toán sinh dữ liệu (vì cần phá vỡ tính đối xứng). \\
        Bảo toàn được khả năng xấp xỉ hàm của mô hình nền. & Hiệu quả giảm mạnh nếu đoán sai tính đối xứng hoặc dùng tập dữ liệu đã căn lề sẵn. \\
        \bottomrule
    \end{tabular}
    \caption[Tóm tắt ưu nhược điểm của geometric machine learning]{Bảng tóm tắt các điểm mạnh và điểm yếu của geometric machine learning}
    \label{tab:uu-nhuoc-diem}
\end{table}

\subsection{Các hướng nghiên cứu tương lai}
\label{subsec:huong-nghien-cuu}

Lĩnh vực geometric machine learning đang phát triển ra nhiều hướng mới. Dựa trên tutorial NeurIPS 2025 \cite{jegelka2025tutorial}, cộng đồng đang tập trung vào một số chủ đề sau.

\subsubsection{Cởi mở hơn với tính đối xứng}

Thay vì gán một loại đối xứng cố định cho toàn bộ mô hình, các nghiên cứu đang tìm cách làm mô hình linh hoạt hơn. Ví dụ, ý tưởng Contextual World Models \cite{gupta2024contextual} thiết kế mô hình tự nhìn vào ngữ cảnh để chọn phép đối xứng cho phù hợp. 

Hướng Any-Subgroup \cite{goel2025anysubgroup} thì đề xuất cố tình bỏ bớt tính đối xứng của toàn hệ thống để đưa bài toán về một nhóm con nhỏ hơn. Điều này giúp giảm tính toán và bắt được các chi tiết cục bộ tốt hơn. Ngoài ra, nhánh Symmetry Discovery đang nghiên cứu thuật toán để mô hình tự tìm ra tính đối xứng trong dữ liệu mà không cần con người phải tự định nghĩa trước.

\subsubsection{Mô hình không phụ thuộc số chiều }

Các mạng nơ-ron thường yêu cầu đầu vào có kích thước cố định (ví dụ ảnh phải là 224x224, hoặc đồ thị có đúng N node). Dựa vào lý thuyết Representation Stability \cite{farb2014representation}, các nhà nghiên cứu đang tạo ra các mô hình có tham số cố định nhưng vẫn nhận được đầu vào có kích thước biến thiên tùy ý. 

Việc này giúp ta có thể huấn luyện mạng GNN trên các đồ thị nhỏ để tiết kiệm chi phí, sau đó mang mạng đó đi dự đoán thẳng trên các đồ thị rất lớn mà kết quả vẫn đúng theo bảo đảm của lý thuyết.

\subsubsection{Kiểm định tính đối xứng}

Thường thì khi code, người ta hay đoán dữ liệu có đối xứng rồi tự thiết kế kiến trúc. Nhánh Symmetry Testing \cite{soleymani2025symmetrytesting} dùng thống kê để kiểm tra xem phân phối dữ liệu có thực sự bảo toàn tính chất sau phép biến đổi hay không. Xu hướng tới là đưa bài kiểm tra này thẳng vào lúc huấn luyện, để mô hình tự động điều chỉnh mức độ áp dụng đối xứng cho khớp với dữ liệu thực tế.

\subsubsection{Học trên không gian trọng số }

Với việc có hàng triệu mô hình mã nguồn mở trên mạng (như Hugging Face), xuất hiện một hướng nghiên cứu là lấy mạng nơ-ron (chính xác là trọng số của nó) làm dữ liệu đầu vào. Trọng số của các lớp ẩn có nhiều phép đối xứng \cite{zhao2026symmetry}, ví dụ đổi chỗ các nơ-ron trong mạng MLP thì kết quả dự đoán cuối cùng không hề thay đổi. 

Các mạng như Neural Functional Transformers \cite{zhou2024neural} hay các kiến trúc equivariant cho trọng số \cite{navon2023equivariant} được tạo ra để cho máy tính đọc hiểu, phân loại và gộp (merge) các mô hình lại với nhau nhanh hơn rất nhiều so với cách thông thường.

\subsubsection{Topo học và độ cong}

Người ta bắt đầu dùng toán học sâu hơn lý thuyết nhóm. Topological Deep Learning \cite{eitan2024topological} dùng phức đơn hình để nhìn được cấu trúc vĩ mô của đồ thị (số lỗ thủng, số liên thông) mà mạng GNN bình thường hay bỏ sót. 

Ngoài ra, việc dùng độ cong không gian (như Ricci curvature) giúp đánh giá xem bài toán học máy dễ hay khó, từ đó tìm ra cách tối ưu mô hình tốt hơn \cite{jegelka2025tutorial}.

\begin{figure}[H]
    \centering
    \begin{tikzpicture}[
        node distance=1.6cm,
        box/.style={draw, rounded corners, align=center, minimum width=4.2cm, minimum height=1cm, font=\small},
        arrow/.style={->, thick}
    ]
        \node[box, fill=green!12] (flex) {Đối xứng linh hoạt\\thích nghi theo ngữ cảnh};
        \node[box, fill=blue!10, right of=flex, xshift=4.5cm] (anydim) {Mô hình không phụ thuộc\\số chiều};
        \node[box, fill=orange!12, below of=flex, yshift=-0.6cm] (test) {Symmetry testing\\tích hợp vào thiết kế};
        \node[box, fill=purple!10, below of=anydim, yshift=-0.6cm] (ws) {Học không gian trọng số\\xử lý Neural Symmetries};
        \node[box, fill=gray!12, below of=test, yshift=-0.6cm] (topo) {Topological DL\\khai thác cấu trúc vĩ mô};
        \node[box, fill=red!10, below of=ws, yshift=-0.6cm] (curv) {Hình học và độ cong\\đánh giá độ khó bài toán};

        \draw[arrow] (flex) -- (test);
        \draw[arrow] (anydim) -- (ws);
        \draw[arrow] (test) -- (topo);
        \draw[arrow] (ws) -- (curv);
    \end{tikzpicture}
    \caption[Sáu hướng nghiên cứu mở trong Geometric Machine Learning]{Các hướng nghiên cứu mới đang được phát triển trong lĩnh vực. Nguồn: Tổng hợp từ \cite{jegelka2025tutorial}}
    \label{fig:open-directions}
\end{figure}

\subsection{Kết luận}
\label{subsec:ket-luan}

Đồ án này bắt đầu bằng thắc mắc: Tại sao mạng CNN dùng cho xử lý ảnh, GNN dùng cho đồ thị hay PointNet dùng cho đám mây điểm lại hiệu quả như vậy? Geometric machine learning trả lời rằng: Các kiến trúc này đã được nhúng sẵn tính đối xứng phù hợp với dữ liệu. Khi mô hình biết hai ảnh thực ra là một do phép biến đổi góc xoay, nó sẽ không lãng phí tài nguyên để học lại bức ảnh đó.

Dù dùng phương pháp can thiệp trên dữ liệu hay xây dựng kiến trúc từ đầu, cách nào cũng có ưu nhược điểm riêng. Data augmentation thì dễ setup và tận dụng được mô hình có sẵn, trong khi các kiến trúc như group convolution thì khó code hơn nhưng lại đảm bảo đúng về mặt toán học.

Những lợi ích này đã được chứng minh rõ qua lý thuyết độ phức tạp mẫu \cite{tahmasebi2023exact}. Dùng tính đối xứng thật sự làm giảm độ phức tạp của bài toán, giúp mô hình cần ít dữ liệu hơn rất nhiều. Vấn đề tốn thời gian tính toán khi phải trung bình qua nguyên một nhóm lớn (group averaging) cũng đã được giảm bớt nhờ các lý thuyết như đồ thị giãn nở \cite{soleymani2025icml}. 

Tất nhiên, geometric ML cũng có nhiều lỗi sai hay gặp. Việc phải tốn thời gian lập trình toán học, khó làm các bài toán sinh dữ liệu, hoặc dễ dùng sai kiến trúc trên tập dữ liệu đã bị căn lề sẵn là những lưu ý quan trọng. Nhưng với các hướng đi mới như học trực tiếp trên trọng số hay kết hợp với topo học, lĩnh vực này chắc chắn sẽ còn phát triển mạnh.

\begin{infobox}{Bài học cốt lõi}
Việc thêm tính chất không gian của dữ liệu vào mạng nơ-ron giúp chúng ta hiểu rằng: mô hình học máy tốt không chỉ cần biết học được gì từ dữ liệu, mà còn cần biết những thuộc tính gì của dữ liệu luôn luôn được giữ nguyên, bất kể chúng ta có biến đổi nó như thế nào.
\end{infobox}

\clearpage
